\section{Class exclusions}\label{sec:classes}

The results of this section exclude classes of machines defined by their head motion. They use the certificate method directly and
reach machines that Theorems~X and~S do not: schedules with macroscopic crossings that are nevertheless excluded, and adaptive machines
with no heavy trace.

\subsection{Boundedly many macroscopic phases}

A walk is \emph{weakly monotone} on an interval of steps if it never moves left there, or never moves right (pauses allowed). It is
\emph{$A$-coarsely monotone} if it never moves more than $A$ cells against a fixed direction: $x(t')\ge x(t)-A$ for all $t\le t'$ in the
interval, or $x(t')\le x(t)+A$ for all such $t\le t'$ (\lean{CoarseMonoOnInterval}).

\begin{definition}[Bounded macro-phases]\label{def:BM}
\formal{MacroBankCrux.lean}{WeakMonoOnInterval, BoundedMacroBank}
Let $M$ be a machine, $C\ge1$, $M_{\mathrm{ph}}\ge1$, $B$ a set of work heads (the \emph{bank}) and $S:\N\to\N$ with $S(d)/L\to0$. $M$ has
\emph{bounded macro-phases} with these parameters (\lean{BoundedMacroBank}) if for all large $d$, whenever $M$ is oblivious at length $L$
with halting time $T\le C(L+1)$, the interval $\Range T$ splits into $M_{\mathrm{ph}}$ consecutive intervals on each of which the input
head and every bank head are weakly monotone, and every other work head visits at most $S(d)$ cells before~$T$.
\end{definition}

The heads outside the bank are unrestricted inside their $o(L)$ workspaces: any number of reversals, any nesting, any coding.

\begin{theorem}[Theorem~B, bounded macro-phases]\label{thm:BM}
\formal{SuffixMassCrux.lean}{boundedMacroBank_has_crux, boundedMacroBank_excludes, suffixMacroBank_has_crux, suffixMacroBank_excludes}
If $M$ has bounded macro-phases with bank $B$ and $M_{\mathrm{ph}}$ phases, then $M$ has crux rate $1/(4(\abs B+1)^{M_{\mathrm{ph}}-1})$ at every
depth, and for all large $d$ there is no $T$ with $\mathrm{ObliviousLinTime}(M,d,C,T)$. The same holds, with crux rate
$\delta_0/(4(\abs B+1)^{M_{\mathrm{ph}}-1})$, if $\delta_0>0$ and the conditions of \cref{def:BM} hold only on a final interval of steps that
contains every append, with at least $\delta_0d$ input rows not read before that interval; the motion before it is unrestricted.
\end{theorem}

\begin{outline}
We prove the suffix form; the whole-run form is its case $\delta_0=1$, with the whole run as the final interval
(\lean{boundedSuffixMacroBank_of_boundedMacroBank}). Dependency paths decrease time, so a path from an output to a read of an untouched
row stays inside the final interval (\lean{getLastD_le_of_mem_depPath}). Let $K_0=(\abs B+1)^{M_{\mathrm{ph}}-1}$ and cut workspace revisits
in the final interval longer than $\lceil\sqrt{S(T+1)}\rceil$; there are $o(L)$ of them. Inside a phase a bank revisit goes back by one step;
between phases, bank revisits are strictly monotone partial injections of fixed sign. So after deleting $o_D(L)$ exceptional outputs,
every short dependency path from an output has a \emph{signature}, the decreasing sequence of phases it enters through bank heads, and at
most $K_0$ signatures start in a given phase. Thresholds on input addresses pull back along signatures to thresholds in time. Split the
untouched input rows into $2K_0$ buckets: each unsplit output row forbids at most $K_0$ buckets, so some bucket is missed by half of the
output rows, and the rectangle of that bucket and those output rows, minus the exceptions, is independent in the cone relation; it has
$(\delta_0/(4K_0)-o(1))L$ elements (\lean{suffixMacroBank_independent_exists}).
\end{outline}

For adaptive machines the phases and the workspaces may depend on the input, and the walks may jitter.

\begin{theorem}[Theorem~J, jittered phases, any adaptivity]\label{thm:J}
\formal{AdaptiveSuffixJitterBM.lean}{native_machine_preserves_suffixJitterBM, nativeGapBound_of_adaptiveJitterBM, adaptiveJitterBM_excludes_transposesAt, adaptiveJitterBM_not_linearTimeTransposer, adaptiveSuffixJitterBM_excludes_transposesAt, adaptiveSuffixJitterBM_not_linearTimeTransposer}
Let $M$ be any machine, $c>0$, $M_{\mathrm{ph}}\ge1$, $B$ a set of work heads, and $A(d),S(d)=o(L)$. Suppose that for all large $d$ every run of $M$ on an input of length $L$ that halts
within $c(L+1)$ steps splits into $M_{\mathrm{ph}}$ intervals on each of which the input head and the bank heads are $A(d)$-coarsely
monotone, while every other work head visits at most $S(d)$ cells; the intervals and cells may depend on the input. Then $M$ does not
transpose within time $c$ at all large $d$; if such $M_{\mathrm{ph}}$, $B$, $A$ and $S$ exist for every $c\ge1$, then $M$ is not a linear-time transposer. The same holds if $\delta_0>0$ and the conditions hold only
on a final interval containing every append with at least $\delta_0d$ rows not entered by the input head before it.
\end{theorem}

\begin{outline}
The bit simulator of the fifteenth paper preserves the shape, also in its suffix form, with the jitter and the workspace enlarged by
$O(s)$ (\lean{native_machine_preserves_suffixJitterBM}). On each frozen trace of the simulator, the proof of Theorem~B goes through with
jitter: the revisits that straddle a pulled-back cut on a bank head occupy at most $2A+1$ cells (\lean{hop_upper_injOn}), and input reads
are pulled back up to a band of width~$A$. This gives every storage code on the frozen cone relations of every simulator rate at most
$1-1/(4K_0)+o(1)$ at the sizes at which $M$ transposes within time~$c$ (\lean{nativeGapBound_of_adaptiveJitterBM}), which is the hypothesis
of \cref{thm:G} below; so Theorem~G excludes~$M$.
\end{outline}

At a fixed number of phases these classes, in their whole-run forms, also yield to pass counting at their own time bounds, by
\cref{thm:coarsepass}, which extends the grouped pass bound of the first paper~\cite[Theorems~5.8 and~5.12]{FirstPaper2026} to coarse
passes with small workspaces. Moreover, the same theorem excludes these shapes, imposed on whole runs of any length, at any running time
for oblivious machines and at any polynomial running time for adaptive ones, even with jitter and workspaces up to a small fixed fraction
of $L$ (\cref{cor:boundedpasses}). What Theorems~B and~J add is an explicit native coding
deficit, and the suffix forms, which pass counting does not reach: an unrestricted prefix may read $(1-\delta_0)d$ rows, and conditioning on the configuration at the start of the suffix
leaves classes of inputs of co-entropy up to $(1-\delta_0)L$, for which the published robust pass bounds~\cite{FirstPaper2026,PassesMemo2026}
say nothing.

\subsection{A native gap suffices}

\paragraph{Roots.} The hypotheses of \cref{thm:G} and \cref{cor:G-uniform,cor:G-low} below are those of the fourteenth paper's native
bound~\cite[Corollary~L]{BoundedStack2026} with a general deficit (\lean{NativeGapBound}). They concern the cone relation of a run
\emph{with arbitrary roots}: for each output position, the root may be any step before the halting step at which the output head is on
that position's cell, and the rate bound is required for every such choice of roots. The append steps are one such choice; whether the
results hold when the bound is required only for them is not proved.

\begin{theorem}[Theorem~G, native gap criterion]\label{thm:G}
\formal{NativeGapCriterion.lean}{nativeGapBound_half_iff, native_gap_excludes_transposesAt, native_gap_excludes_linear}
Let $M$ be a machine and $c>0$. Suppose that there is $\rho>0$ such that, for every $s\ge1$ and every depth $D$, every binary storage
code on the cone relation with arbitrary roots, after a genie of full depth $D$, of every run of the simulator $\Ns$ of $M$~\cite[Theorem~S]{HighRate2026}
that halts within $(11c+1)(L+1)$ steps on an input of length $L$ (entries of $b(d)$ bits) has rate at most $1-\rho+o(1)$, where $d$ ranges
over the sizes at which $M$ transposes within time~$c$. Then $M$ transposes within time $c$ at only finitely many $d$. If such a $\rho$
exists for every machine $M$ and every $c>0$, then (T) fails.
\end{theorem}

\begin{outline}
Choose $s$ so that, at all large $d$ at which $M$ transposes within time $c$, $\Ns$ transposes within time $11c+1$ with at most
$2^{(\rho/2)L}$ traces~\cite[Theorem~S]{HighRate2026}. Some trace of $\Ns$ carries $2^{(1-\rho/2)L}$ inputs; after the genie, of cost
$\frac14\rho L$, its fibre splits into storage codes on its cone relation (\lean{MergeDepth.blockMachine_isStorageCode_fibre}), one of
which has rate at least $1-\frac34\rho-o(1)$, contradicting the hypothesis. If (T) held, some machine $M$ would transpose in time
$c(L+1)$ at infinitely many sizes. At $\rho=\frac12$ the hypothesis is the native half bound of the fourteenth
paper~\cite[Corollary~L]{BoundedStack2026} on the simulators of~$M$.
\end{outline}

\begin{corollary}[Theorem~G, uniform form]\label{cor:G-uniform}
\formal{NativeGapCriterion.lean}{native_gap_exists_exclude_linear, native_gap_tests_exclude_linear}
Suppose that for every $k$ and $C\ge1$ there is $\rho(k,C)>0$ such that, for every machine with $k$ work tapes and every depth $D$, every
binary storage code on the cone relation with arbitrary roots, after a genie of full depth $D$, of every run that halts within $C(L+1)$ steps on an input of
length $L$ (entries of $b(d)$ bits) has rate at most $1-\rho(k,C)+o(1)$. Then (T) fails.
\end{corollary}

\begin{outline}
The simulators of a machine $M$ with $k$ work tapes have $k$ work tapes (\lean{native_machine_k}), so for every $c>0$ the deficit
$\rho(k,11c+1)$ satisfies the hypothesis of \cref{thm:G} for $M$ and $c$. At $\rho=\frac12$ the hypothesis is the native half bound of the
fourteenth paper~\cite[Corollary~L]{BoundedStack2026}, assumed uniformly over the machines with $k$ work tapes and over all depths, as in
the eighth paper~\cite[\S6.2]{EighthPaper2026}.
\end{outline}

The fifteenth paper reaches the same conclusion from the half bound on chosen simulators~\cite{HighRate2026}; Theorem~G accepts any
fixed deficit, chosen for each machine and time bound separately and needed only on that machine's simulators. In the uniform form of
\cref{cor:G-uniform} the hypothesis is at least as strong as a uniform crux; its value is that the deficit may be tiny, which is the
form that Theorems~B and~J produce, but only for their classes.

For a single machine of low adaptivity the same counting needs no simulator: a deficit above the entropy of the trace suffices.

\begin{corollary}[A native gap above the trace entropy]\label{cor:G-low}
\formal{NativeGapCriterion.lean}{fSecond_native_gap, native_gap_exclusion_of_trace_support}
Let $M$ be a machine with $k$ work tapes, let $c>0$ and $\eta<\rho$, let $\beta_1$ be the block width of $M$ at scale~$1$
(\lean{MergeDepth.blockWidth}), and put $D=\mathrm{genieDepth}\bigl(k,\beta_1,\max(c,1),(\rho-\eta)/2\bigr)$. Suppose that every binary storage code on
the cone relation with arbitrary roots, after a genie of full depth $D$, of every run of $M$ that halts within $c(L+1)$ steps has rate at most
$1-\rho+o(1)$. Then at all large $d$ at which $M$ transposes within time $c$, the skeleton of $M$ at scale~$1$, and hence its trace, has
entropy greater than $\eta L$ on a uniformly random input (\lean{TapeStorageRates.skeletonOfTrace_eq}, \lean{MergeDepth.entropy_comp_le}).
Moreover, if a machine $N$ satisfies the same hypothesis, with its own $k$, $\beta_1$ and $D$, and for all large $d$ every $d$ at which $M$
transposes within time $c$ is one at which $N$ transposes within time $c$ with at most $2^{\eta L}$ distinct traces, then $M$ transposes
within time $c$ at only finitely many~$d$.
\end{corollary}

\begin{outline}
If the trace entropy were at most $\eta L$, some trace would be followed by at least $2^{(1-\eta)L}$ inputs. After the genie of cost
$\frac12(\rho-\eta)L$, its fibre splits into storage codes on the cone relation, one of which has rate at least
$1-\eta-\frac12(\rho-\eta)>1-\rho$, contradicting the hypothesis. This is the argument of \cref{thm:G} with $M$ in place of the
simulator; at $\rho=\frac12$ it is the fourteenth paper's low-adaptivity corollary~\cite[Corollary~L]{BoundedStack2026}.
\end{outline}

\subsection{Few scales and no free head}

For a work head, a \emph{family} is the set of its revisits on one side that share their \emph{turning point}, the first step inside the
revisit at which the head is farthest from the cell; its \emph{height} is the distance from the turning point to the family's outermost
cell, the cell of its longest revisit, as in the eighth paper~\cite{EighthPaper2026}. The head has \emph{owner load} at most $\ell$ if
at every step, on each side, the revisits open at that step have at most $\ell$ distinct turning points. Bounded cell entry $V$ gives owner
load at most $2V-1$.

\begin{theorem}[Theorem~Q, scale-gap pigeonhole]\label{thm:Q}
\formal{ZeroFreeHeadQWindow.lean}{zeroFreeHead_qWindow_excludes, zeroFreeHead_entries_qWindow_excludes, zeroFreeHead_qWindow_oblivious_excludes}
Let $c,\beta$ be reals and $k,\ell,q,h$ natural numbers. There are $\eps>0$ and $d_0$ such that for every $\eta<\eps$ and $d\ge d_0$ no machine with at
most $k$ work tapes and value width at most $\beta$ transposes at size $d$ within time $c$ with an $\eta$-heavy trace on which every work
head has owner load at most $\ell$ and the binary logarithms of the heights of all families lie in at most $q$ windows of width~$h$. The
input head is unrestricted. For oblivious machines the statement holds at every $d\ge d_0$ with $\eta=0$.
\end{theorem}

\begin{outline}
Owner load $\ell$ colours the families with $N=2k\max(\ell,1)$ classes of non-overlapping spans, so their heights sum to less than $NcL$. Choose
thresholds $Q_0<Q_1<\dots<Q_q$, each so large compared with the last that a window can be \emph{bad}, that is contain the height of a
family longer than the current threshold, for at most one index. With $q$ windows and $q+1$ indices some index is bad for no window, so
only boundedly many families are large, and the eleventh paper's low-adaptivity theorem for boundedly many large
families~\cite{MergeDepth2026} gives the contradiction. Its constant comes from a compactness argument and is not
explicit~\cite{SourceCovers2026}.
\end{outline}

Theorem~Q is the case without a free work head of the eighth paper's bounded-load, boundedly-many-scales class, which is
conditional on the crux there~\cite{EighthPaper2026}; here it is unconditional, with any input motion and low adaptivity. With
$\Theta(\log d)$ scales instead of boundedly many, the question is open already for two work tapes.
